\documentclass[10pt,a4paper]{amsart}
\usepackage{amsmath,amssymb,amsthm,graphicx}
% Retained minimal preamble for the Stacks Project chapter more-algebra.tex (More on Algebra).
% The chapter file pinned-original-source.tex carries only \input{preamble} on its first line, so no
% preamble text lives inside the chapter itself.  The Stacks Project uses one single shared preamble for
% every chapter; the snapshot stages it identically as sources/stacks-chow/preamble.tex and
% sources/stacks-etale-cohomology/preamble.tex.  The block below is copied mechanically, line for line,
% from sources/stacks-chow/preamble.tex lines 14-18 (the Xy-pic setup) and 154-217 (the native theorem
% declarations together with the mathematical macros).  Nothing was added, removed or redefined; the
% wrapper class and the amsmath/amssymb/amsthm/graphicx base packages are supplied by the assembler.
% Chapter: more-algebra.tex -- More on Algebra.
% For commutative diagrams we use Xy-pic
\usepackage[all]{xy}

% We use 2cell for 2-commutative diagrams.
\xyoption{2cell}

% Theorem environments.
%
\theoremstyle{plain}
\newtheorem{theorem}[subsection]{Theorem}
\newtheorem{proposition}[subsection]{Proposition}
\newtheorem{lemma}[subsection]{Lemma}

\theoremstyle{definition}
\newtheorem{definition}[subsection]{Definition}
\newtheorem{example}[subsection]{Example}
\newtheorem{exercise}[subsection]{Exercise}
\newtheorem{situation}[subsection]{Situation}

\theoremstyle{remark}
\newtheorem{remark}[subsection]{Remark}
\newtheorem{remarks}[subsection]{Remarks}

\numberwithin{equation}{subsection}

% Macros
%
\def\lim{\mathop{\mathrm{lim}}\nolimits}
\def\colim{\mathop{\mathrm{colim}}\nolimits}
\def\Spec{\mathop{\mathrm{Spec}}}
\def\Hom{\mathop{\mathrm{Hom}}\nolimits}
\def\Ext{\mathop{\mathrm{Ext}}\nolimits}
\def\SheafHom{\mathop{\mathcal{H}\!\mathit{om}}\nolimits}
\def\SheafExt{\mathop{\mathcal{E}\!\mathit{xt}}\nolimits}
\def\Sch{\mathit{Sch}}
\def\Mor{\mathop{\mathrm{Mor}}\nolimits}
\def\Ob{\mathop{\mathrm{Ob}}\nolimits}
\def\Sh{\mathop{\mathit{Sh}}\nolimits}
\def\NL{\mathop{N\!L}\nolimits}
\def\CH{\mathop{\mathrm{CH}}\nolimits}
\def\proetale{{pro\text{-}\acute{e}tale}}
\def\etale{{\acute{e}tale}}
\def\QCoh{\mathit{QCoh}}
\def\Ker{\mathop{\mathrm{Ker}}}
\def\Im{\mathop{\mathrm{Im}}}
\def\Coker{\mathop{\mathrm{Coker}}}
\def\Coim{\mathop{\mathrm{Coim}}}

% Boxtimes
%
\DeclareMathSymbol{\boxtimes}{\mathbin}{AMSa}{"02}

%
% Macros for moduli stacks/spaces
%
\def\QCohstack{\mathcal{QC}\!\mathit{oh}}
\def\Cohstack{\mathcal{C}\!\mathit{oh}}
\def\Spacesstack{\mathcal{S}\!\mathit{paces}}
\def\Quotfunctor{\mathrm{Quot}}
\def\Hilbfunctor{\mathrm{Hilb}}
\def\Curvesstack{\mathcal{C}\!\mathit{urves}}
\def\Polarizedstack{\mathcal{P}\!\mathit{olarized}}
\def\Complexesstack{\mathcal{C}\!\mathit{omplexes}}
% \Pic is the operator that assigns to X its picard group, usage \Pic(X)
% \Picardstack_{X/B} denotes the Picard stack of X over B
% \Picardfunctor_{X/B} denotes the Picard functor of X over B
\def\Pic{\mathop{\mathrm{Pic}}\nolimits}
\def\Picardstack{\mathcal{P}\!\mathit{ic}}
\def\Picardfunctor{\mathrm{Pic}}
\def\Deformationcategory{\mathcal{D}\!\mathit{ef}}

% --- appended 2026-10-12 (Chief): the project own definitions of the reference/slogan/history
% display environments, copied mechanically from the shared Stacks preamble lines 8-12
% (\usepackage{verbatim} plus the three \newenvironment lines, which the chapter preamble keeps
% verbatim from the project preamble). Nothing added, removed or redefined: these environments
% carry no mathematical content and the project itself typesets them as comments. Without these
% three lines a counted statement whose own text uses them fails with Environment slogan undefined.
\usepackage{verbatim}
\newenvironment{reference}{\comment}{\endcomment}
\newenvironment{slogan}{\comment}{\endcomment}
\newenvironment{history}{\comment}{\endcomment}
\usepackage{cleveref}
\crefname{originallemma}{original lemma}{original lemmas}
\Crefname{originallemma}{Original lemma}{Original lemmas}
\makeatletter
\newlabel{algebra-lemma-tensor-products-commute-with-limits}{{00DD}{source}{}{sourcepointer-1}{}}
\newlabel{algebra-lemma-tensor-products-commute-with-limits@cref}{{[originallemma][0][]00DD}{[1][0][]source}{}{}{}}
\makeatother
\crefname{originallemma}{original lemma}{original lemmas}
\Crefname{originallemma}{Original lemma}{Original lemmas}
\makeatletter
\newlabel{algebra-lemma-hathat-finitely-generated}{{05GG}{source}{}{sourcepointer-2}{}}
\newlabel{algebra-lemma-hathat-finitely-generated@cref}{{[originallemma][0][]05GG}{[1][0][]source}{}{}{}}
\makeatother
\title{A power torsion module is unchanged by base change along a map that is an isomorphism on the quotients by the powers of the ideal}
\author{The Stacks Project Authors}
\date{stacks-project:more-algebra.tex:local-snapshot-20261012; GNU FDL 1.2 or later}
\begin{document}
\maketitle
\setcounter{subsection}{1}
\noindent\textit{Editorial note.} All delivered mathematical text is the author's own, copied from the pinned Stacks Project chapter file more-algebra.tex; no statement or proof step was rewritten, reordered or removed. The counted statement is the labeled lemma environment at source lines 24172-24183 (canonical tag 0BNK) and the counted proof is the proof environment at source lines 24185-24198, which immediately follows it and is its own proof.\par\medskip
\noindent\textit{Editorial note.} The counted statement contains the author's reference environment recording that the statement is a slight generalization of a result of Beauville and Laszlo; the retained preamble defines that environment as comment-type, so its body and the citation inside it are discarded when the document is typeset, which is why the card needs no bibliography.\par\medskip
\noindent\textit{Editorial note.} The two references are to results of another chapter on commutation of tensor products with limits and on finitely generated ideals; the kinds recorded are the type words appearing in the labels themselves. Both are navigation stubs with their roles recorded in the reasons.\par\medskip
\noindent\textit{Editorial note.} The preamble span runs to line 90, the current end of the retained preamble, which is exactly what makes the reference environment a comment; the excerpt serves only as preamble, never as a span coordinate.\par\medskip
\noindent\textit{Editorial note.} The source URL is the source-level Stacks Project URL, as required when more than one card is produced from one source; the per-claim tag is carried by claim\_key and by this note.\par\medskip
\noindent\textit{Allowed standard prerequisites (editorial).}\begin{itemize}
\item A module annihilated by a power of the ideal is a module over the quotient, and tensoring it over the quotient against the base changed quotient gives back the module when the quotients agree. used at proof lines 24185-24190. Establishes the case of a module annihilated by a fixed power through the displayed chain of isomorphisms.
\item A power torsion module is the union of the submodules annihilated by the successive powers of the ideal, and tensor products commute with direct limits. used at proof lines 24190-24196. Passes from the fixed powers to an arbitrary power torsion module.
\item The completion of the ring with respect to a finitely generated ideal is a special case of the base change setting, namely the base change along the completion map. used at proof lines 24196-24198. Deduces the final assertion of the statement from the general one.
\end{itemize}\medskip
\ifdefined\hypertarget\hypertarget{sourcepointer-1}{}\fi
\noindent\textit{Source reference.} Original lemma 00DD: sources/stacks-chow/tags.txt, source lines 493--493. The reference is invoked at proof lines 24190-24193 to pass from the modules annihilated by a fixed power to an arbitrary power torsion module; the chain of identifications for the fixed power is inside the counted span.\par\medskip
\ifdefined\hypertarget\hypertarget{sourcepointer-2}{}\fi
\noindent\textit{Source reference.} Original lemma 05GG: sources/stacks-chow/tags.txt, source lines 7060--7060. The reference is cited at proof lines 24196-24198 for the final assertion, identifying the completion of a finitely generated ideal as an instance of the general setting; it is a pointer and the general statement is proved inside the counted span.\par\medskip
\begin{lemma}
\label{lemma-torsion-completion}
\begin{reference}
Slight generalization of \cite[Lemme~1]{Beauville-Laszlo}.
\end{reference}
Let $R \to R'$ be a ring map. Let $I \subset R$ be an ideal such that
$R/I^n \to R'/I^nR'$ is an isomorphism for $n > 0$.
For any $I$-power torsion $R$-module $M$ the map $M \to M \otimes_R R'$
is an isomorphism.
For example, if $I$ is finitely generated and $R^\wedge$ is the completion
of $R$ with respect to $I$, then we have $M \cong M \otimes_R R^\wedge$.
\end{lemma}

\begin{proof}
If $M$ is annihilated by $I^n$, then
$$
M \otimes_R R' \cong
M \otimes_{R/I^n} R'/I^n R' \cong
M \otimes_{R/I^n} R/I^n \cong M.
$$
If $M$ is $I$-power torsion, then $M = \bigcup M[I^n]$.
Since tensor products commute with direct limits
(Algebra, Lemma \ref{algebra-lemma-tensor-products-commute-with-limits}),
we obtain the desired isomorphism. The last statement is a special
case of the first statement by Algebra, Lemma
\ref{algebra-lemma-hathat-finitely-generated}.
\end{proof}

\end{document}
